% Options for packages loaded elsewhere
\PassOptionsToPackage{unicode}{hyperref}
\PassOptionsToPackage{hyphens}{url}
\documentclass[
]{article}
\usepackage{xcolor}
\usepackage[margin=1in]{geometry}
\usepackage{amsmath,amssymb}
\setcounter{secnumdepth}{-\maxdimen} % remove section numbering
\usepackage{iftex}
\ifPDFTeX
  \usepackage[T1]{fontenc}
  \usepackage[utf8]{inputenc}
  \usepackage{textcomp} % provide euro and other symbols
\else % if luatex or xetex
  \usepackage{unicode-math} % this also loads fontspec
  \defaultfontfeatures{Scale=MatchLowercase}
  \defaultfontfeatures[\rmfamily]{Ligatures=TeX,Scale=1}
\fi
\usepackage{lmodern}
\ifPDFTeX\else
  % xetex/luatex font selection
\fi
% Use upquote if available, for straight quotes in verbatim environments
\IfFileExists{upquote.sty}{\usepackage{upquote}}{}
\IfFileExists{microtype.sty}{% use microtype if available
  \usepackage[]{microtype}
  \UseMicrotypeSet[protrusion]{basicmath} % disable protrusion for tt fonts
}{}
\makeatletter
\@ifundefined{KOMAClassName}{% if non-KOMA class
  \IfFileExists{parskip.sty}{%
    \usepackage{parskip}
  }{% else
    \setlength{\parindent}{0pt}
    \setlength{\parskip}{6pt plus 2pt minus 1pt}}
}{% if KOMA class
  \KOMAoptions{parskip=half}}
\makeatother
\usepackage{color}
\usepackage{fancyvrb}
\newcommand{\VerbBar}{|}
\newcommand{\VERB}{\Verb[commandchars=\\\{\}]}
\DefineVerbatimEnvironment{Highlighting}{Verbatim}{commandchars=\\\{\}}
% Add ',fontsize=\small' for more characters per line
\usepackage{framed}
\definecolor{shadecolor}{RGB}{248,248,248}
\newenvironment{Shaded}{\begin{snugshade}}{\end{snugshade}}
\newcommand{\AlertTok}[1]{\textcolor[rgb]{0.94,0.16,0.16}{#1}}
\newcommand{\AnnotationTok}[1]{\textcolor[rgb]{0.56,0.35,0.01}{\textbf{\textit{#1}}}}
\newcommand{\AttributeTok}[1]{\textcolor[rgb]{0.13,0.29,0.53}{#1}}
\newcommand{\BaseNTok}[1]{\textcolor[rgb]{0.00,0.00,0.81}{#1}}
\newcommand{\BuiltInTok}[1]{#1}
\newcommand{\CharTok}[1]{\textcolor[rgb]{0.31,0.60,0.02}{#1}}
\newcommand{\CommentTok}[1]{\textcolor[rgb]{0.56,0.35,0.01}{\textit{#1}}}
\newcommand{\CommentVarTok}[1]{\textcolor[rgb]{0.56,0.35,0.01}{\textbf{\textit{#1}}}}
\newcommand{\ConstantTok}[1]{\textcolor[rgb]{0.56,0.35,0.01}{#1}}
\newcommand{\ControlFlowTok}[1]{\textcolor[rgb]{0.13,0.29,0.53}{\textbf{#1}}}
\newcommand{\DataTypeTok}[1]{\textcolor[rgb]{0.13,0.29,0.53}{#1}}
\newcommand{\DecValTok}[1]{\textcolor[rgb]{0.00,0.00,0.81}{#1}}
\newcommand{\DocumentationTok}[1]{\textcolor[rgb]{0.56,0.35,0.01}{\textbf{\textit{#1}}}}
\newcommand{\ErrorTok}[1]{\textcolor[rgb]{0.64,0.00,0.00}{\textbf{#1}}}
\newcommand{\ExtensionTok}[1]{#1}
\newcommand{\FloatTok}[1]{\textcolor[rgb]{0.00,0.00,0.81}{#1}}
\newcommand{\FunctionTok}[1]{\textcolor[rgb]{0.13,0.29,0.53}{\textbf{#1}}}
\newcommand{\ImportTok}[1]{#1}
\newcommand{\InformationTok}[1]{\textcolor[rgb]{0.56,0.35,0.01}{\textbf{\textit{#1}}}}
\newcommand{\KeywordTok}[1]{\textcolor[rgb]{0.13,0.29,0.53}{\textbf{#1}}}
\newcommand{\NormalTok}[1]{#1}
\newcommand{\OperatorTok}[1]{\textcolor[rgb]{0.81,0.36,0.00}{\textbf{#1}}}
\newcommand{\OtherTok}[1]{\textcolor[rgb]{0.56,0.35,0.01}{#1}}
\newcommand{\PreprocessorTok}[1]{\textcolor[rgb]{0.56,0.35,0.01}{\textit{#1}}}
\newcommand{\RegionMarkerTok}[1]{#1}
\newcommand{\SpecialCharTok}[1]{\textcolor[rgb]{0.81,0.36,0.00}{\textbf{#1}}}
\newcommand{\SpecialStringTok}[1]{\textcolor[rgb]{0.31,0.60,0.02}{#1}}
\newcommand{\StringTok}[1]{\textcolor[rgb]{0.31,0.60,0.02}{#1}}
\newcommand{\VariableTok}[1]{\textcolor[rgb]{0.00,0.00,0.00}{#1}}
\newcommand{\VerbatimStringTok}[1]{\textcolor[rgb]{0.31,0.60,0.02}{#1}}
\newcommand{\WarningTok}[1]{\textcolor[rgb]{0.56,0.35,0.01}{\textbf{\textit{#1}}}}
\usepackage{graphicx}
\makeatletter
\newsavebox\pandoc@box
\newcommand*\pandocbounded[1]{% scales image to fit in text height/width
  \sbox\pandoc@box{#1}%
  \Gscale@div\@tempa{\textheight}{\dimexpr\ht\pandoc@box+\dp\pandoc@box\relax}%
  \Gscale@div\@tempb{\linewidth}{\wd\pandoc@box}%
  \ifdim\@tempb\p@<\@tempa\p@\let\@tempa\@tempb\fi% select the smaller of both
  \ifdim\@tempa\p@<\p@\scalebox{\@tempa}{\usebox\pandoc@box}%
  \else\usebox{\pandoc@box}%
  \fi%
}
% Set default figure placement to htbp
\def\fps@figure{htbp}
\makeatother
\setlength{\emergencystretch}{3em} % prevent overfull lines
\providecommand{\tightlist}{%
  \setlength{\itemsep}{0pt}\setlength{\parskip}{0pt}}
\usepackage{bookmark}
\IfFileExists{xurl.sty}{\usepackage{xurl}}{} % add URL line breaks if available
\urlstyle{same}
\hypersetup{
  pdftitle={scRNA-seq Analysis of PBMC 3k Dataset},
  pdfauthor={Dhvani Solanki},
  hidelinks,
  pdfcreator={LaTeX via pandoc}}

\title{scRNA-seq Analysis of PBMC 3k Dataset}
\author{Dhvani Solanki}
\date{2026-04-16}

\begin{document}
\maketitle

\begin{center}\rule{0.5\linewidth}{0.5pt}\end{center}

title: ``scRNA-seq Analysis of PBMC 3k Dataset'' author: ``Dhwanit''
date: ``2026-04-16'' output: html\_document ---------------------

\subsection{Overview}\label{overview}

This project performs an end-to-end single-cell RNA-seq analysis of the
PBMC 3k dataset using Seurat. The workflow includes quality control,
normalization, dimensionality reduction, clustering, marker
identification, and cell type annotation.

\begin{center}\rule{0.5\linewidth}{0.5pt}\end{center}

\subsection{1. Load Libraries}\label{load-libraries}

\begin{Shaded}
\begin{Highlighting}[]
\FunctionTok{library}\NormalTok{(Seurat)}
\end{Highlighting}
\end{Shaded}

\begin{verbatim}
## Loading required package: SeuratObject
\end{verbatim}

\begin{verbatim}
## Loading required package: sp
\end{verbatim}

\begin{verbatim}
## 'SeuratObject' was built under R 4.5.2 but the current version is
## 4.5.3; it is recomended that you reinstall 'SeuratObject' as the ABI
## for R may have changed
\end{verbatim}

\begin{verbatim}
## 
## Attaching package: 'SeuratObject'
\end{verbatim}

\begin{verbatim}
## The following objects are masked from 'package:base':
## 
##     intersect, t
\end{verbatim}

\begin{Shaded}
\begin{Highlighting}[]
\FunctionTok{library}\NormalTok{(tidyverse)}
\end{Highlighting}
\end{Shaded}

\begin{verbatim}
## -- Attaching core tidyverse packages ------------------------ tidyverse 2.0.0 --
## v dplyr     1.1.4     v readr     2.1.5
## v forcats   1.0.0     v stringr   1.5.1
## v ggplot2   4.0.2     v tibble    3.2.1
## v lubridate 1.9.4     v tidyr     1.3.1
## v purrr     1.0.4
\end{verbatim}

\begin{verbatim}
## -- Conflicts ------------------------------------------ tidyverse_conflicts() --
## x dplyr::filter() masks stats::filter()
## x dplyr::lag()    masks stats::lag()
## i Use the conflicted package (<http://conflicted.r-lib.org/>) to force all conflicts to become errors
\end{verbatim}

\begin{Shaded}
\begin{Highlighting}[]
\FunctionTok{library}\NormalTok{(cowplot)}
\end{Highlighting}
\end{Shaded}

\begin{verbatim}
## 
## Attaching package: 'cowplot'
## 
## The following object is masked from 'package:lubridate':
## 
##     stamp
\end{verbatim}

\begin{Shaded}
\begin{Highlighting}[]
\FunctionTok{library}\NormalTok{(ggplot2)}
\end{Highlighting}
\end{Shaded}

\begin{center}\rule{0.5\linewidth}{0.5pt}\end{center}

\subsection{2. Load Data}\label{load-data}

\begin{Shaded}
\begin{Highlighting}[]
\NormalTok{pbmc.data }\OtherTok{\textless{}{-}} \FunctionTok{Read10X}\NormalTok{(}\AttributeTok{data.dir =} \StringTok{"data/filtered\_gene\_bc\_matrices/hg19/"}\NormalTok{)}

\NormalTok{pbmc }\OtherTok{\textless{}{-}} \FunctionTok{CreateSeuratObject}\NormalTok{(}\AttributeTok{counts =}\NormalTok{ pbmc.data,}
                           \AttributeTok{project =} \StringTok{"pbmc3k"}\NormalTok{,}
                           \AttributeTok{min.cells =} \DecValTok{3}\NormalTok{,}
                           \AttributeTok{min.features =} \DecValTok{200}\NormalTok{)}
\end{Highlighting}
\end{Shaded}

\begin{verbatim}
## Warning: Feature names cannot have underscores ('_'), replacing with dashes
## ('-')
\end{verbatim}

\begin{Shaded}
\begin{Highlighting}[]
\NormalTok{pbmc}
\end{Highlighting}
\end{Shaded}

\begin{verbatim}
## An object of class Seurat 
## 13714 features across 2700 samples within 1 assay 
## Active assay: RNA (13714 features, 0 variable features)
##  1 layer present: counts
\end{verbatim}

\begin{center}\rule{0.5\linewidth}{0.5pt}\end{center}

\subsection{3. Quality Control (QC)}\label{quality-control-qc}

\begin{Shaded}
\begin{Highlighting}[]
\NormalTok{pbmc[[}\StringTok{"percent.mt"}\NormalTok{]] }\OtherTok{\textless{}{-}} \FunctionTok{PercentageFeatureSet}\NormalTok{(pbmc, }\AttributeTok{pattern =} \StringTok{"\^{}MT{-}"}\NormalTok{)}

\FunctionTok{VlnPlot}\NormalTok{(pbmc, }\AttributeTok{features =} \FunctionTok{c}\NormalTok{(}\StringTok{"nFeature\_RNA"}\NormalTok{, }\StringTok{"nCount\_RNA"}\NormalTok{, }\StringTok{"percent.mt"}\NormalTok{), }\AttributeTok{ncol =} \DecValTok{3}\NormalTok{)}
\end{Highlighting}
\end{Shaded}

\begin{verbatim}
## Warning: Default search for "data" layer in "RNA" assay yielded no results;
## utilizing "counts" layer instead.
\end{verbatim}

\pandocbounded{\includegraphics[keepaspectratio]{pbmc_analysis_files/pbmc_analysis_files/figure-latex/unnamed-chunk-3-1.pdf}}

\subsubsection{Filter low-quality cells}\label{filter-low-quality-cells}

\begin{Shaded}
\begin{Highlighting}[]
\NormalTok{pbmc }\OtherTok{\textless{}{-}} \FunctionTok{subset}\NormalTok{(pbmc, }\AttributeTok{subset =}\NormalTok{ nFeature\_RNA }\SpecialCharTok{\textgreater{}} \DecValTok{200} \SpecialCharTok{\&} 
\NormalTok{                              nFeature\_RNA }\SpecialCharTok{\textless{}} \DecValTok{2500} \SpecialCharTok{\&} 
\NormalTok{                              percent.mt }\SpecialCharTok{\textless{}} \DecValTok{5}\NormalTok{)}

\NormalTok{pbmc}
\end{Highlighting}
\end{Shaded}

\begin{verbatim}
## An object of class Seurat 
## 13714 features across 2638 samples within 1 assay 
## Active assay: RNA (13714 features, 0 variable features)
##  1 layer present: counts
\end{verbatim}

\begin{center}\rule{0.5\linewidth}{0.5pt}\end{center}

\subsection{4. Normalization and Feature
Selection}\label{normalization-and-feature-selection}

\begin{Shaded}
\begin{Highlighting}[]
\NormalTok{pbmc }\OtherTok{\textless{}{-}} \FunctionTok{NormalizeData}\NormalTok{(pbmc)}
\end{Highlighting}
\end{Shaded}

\begin{verbatim}
## Normalizing layer: counts
\end{verbatim}

\begin{Shaded}
\begin{Highlighting}[]
\NormalTok{pbmc }\OtherTok{\textless{}{-}} \FunctionTok{FindVariableFeatures}\NormalTok{(pbmc, }\AttributeTok{selection.method =} \StringTok{"vst"}\NormalTok{, }\AttributeTok{nfeatures =} \DecValTok{1500}\NormalTok{)}
\end{Highlighting}
\end{Shaded}

\begin{verbatim}
## Finding variable features for layer counts
\end{verbatim}

\begin{center}\rule{0.5\linewidth}{0.5pt}\end{center}

\subsection{5. Scaling and PCA}\label{scaling-and-pca}

\begin{Shaded}
\begin{Highlighting}[]
\NormalTok{pbmc }\OtherTok{\textless{}{-}} \FunctionTok{ScaleData}\NormalTok{(pbmc)}
\end{Highlighting}
\end{Shaded}

\begin{verbatim}
## Centering and scaling data matrix
\end{verbatim}

\begin{Shaded}
\begin{Highlighting}[]
\NormalTok{pbmc }\OtherTok{\textless{}{-}} \FunctionTok{RunPCA}\NormalTok{(pbmc)}
\end{Highlighting}
\end{Shaded}

\begin{verbatim}
## PC_ 1 
## Positive:  CST3, TYROBP, LST1, AIF1, FTL, FTH1, LYZ, FCN1, S100A9, TYMP 
##     FCER1G, CFD, LGALS1, S100A8, LGALS2, CTSS, SERPINA1, IFITM3, CFP, PSAP 
##     SAT1, IFI30, S100A11, COTL1, NPC2, LGALS3, GSTP1, PYCARD, NCF2, CDA 
## Negative:  LTB, CD2, ACAP1, STK17A, CTSW, GIMAP5, AQP3, CCL5, GZMA, CST7 
##     MAL, ITM2A, HOPX, LDLRAP1, GZMK, ZAP70, ETS1, TNFAIP8, RIC3, SAMD3 
##     LYAR, KLRG1, RORA, NKG7, TRABD2A, PRF1, FAM107B, CCND3, CD79A, SRSF7 
## PC_ 2 
## Positive:  CD79A, MS4A1, TCL1A, HLA-DQA1, HLA-DQB1, HLA-DRA, LINC00926, CD79B, CD74, HLA-DRB1 
##     HLA-DPB1, HLA-DMA, HLA-DRB5, HLA-DPA1, LTB, FCRLA, HVCN1, BLNK, P2RX5, IRF8 
##     IGLL5, ARHGAP24, SMIM14, PPP1R14A, C16orf74, MZB1, RP11-428G5.5, RP5-887A10.1, IL4R, BTK 
## Negative:  NKG7, PRF1, CST7, GZMA, GZMB, FGFBP2, CTSW, GNLY, SPON2, GZMH 
##     CCL4, CCL5, FCGR3A, CLIC3, XCL2, HOPX, SRGN, TTC38, S100A4, IGFBP7 
##     ID2, ANXA1, ACTB, SAMD3, APOBEC3G, KLRG1, PTGDR, ABI3, LYAR, IFITM2 
## PC_ 3 
## Positive:  HLA-DPA1, HLA-DPB1, HLA-DQA1, HLA-DQB1, CD79A, CD79B, CD74, HLA-DRB1, MS4A1, HLA-DRB5 
##     HLA-DRA, TCL1A, LINC00926, HLA-DMA, HVCN1, IRF8, FCRLA, SMIM14, BLNK, LAT2 
##     P2RX5, IGLL5, ARHGAP24, RBM3, CSNK2B, IFITM2, CD1C, UBE2L6, APOBEC3G, LTB 
## Negative:  PPBP, PF4, SDPR, SPARC, GNG11, NRGN, GP9, HIST1H2AC, RGS18, TUBB1 
##     CLU, AP001189.4, ITGA2B, CD9, PTCRA, TMEM40, CA2, ACRBP, MMD, TREML1 
##     SEPT5, RUFY1, TSC22D1, MYL9, MPP1, CMTM5, RP11-367G6.3, LY6G6F, GP1BA, NGFRAP1 
## PC_ 4 
## Positive:  HLA-DQA1, CD79B, CD79A, HLA-DQB1, CD74, MS4A1, HLA-DPB1, HLA-DPA1, HLA-DRB1, HLA-DRA 
##     TCL1A, HLA-DRB5, LINC00926, HLA-DMA, HVCN1, GZMB, FCRLA, FGFBP2, IRF8, BLNK 
##     FCGR3A, PRF1, NKG7, IGLL5, P2RX5, GNLY, SMIM14, CST7, LAT2, SPON2 
## Negative:  S100A6, S100A8, S100A4, FYB, MAL, AQP3, S100A9, TMSB4X, CD2, GIMAP4 
##     RBP7, LGALS2, ANXA1, GIMAP5, S100A12, TRABD2A, FCN1, LYZ, FOLR3, MS4A6A 
##     LDLRAP1, S100A11, AIF1, IL8, PASK, ITM2A, CORO1B, ATP5H, ASGR1, RETN 
## PC_ 5 
## Positive:  LTB, CKB, MS4A7, AQP3, PTGES3, HN1, CORO1B, LILRB2, CD2, GDI2 
##     IFITM2, MAL, ANXA5, ABRACL, CTD-2006K23.1, WARS, PPA1, ATP5C1, VMO1, NAAA 
##     HSP90AA1, FYB, GIMAP4, ARL6IP5, ADA, COTL1, GPR183, PPM1N, ITM2A, TIMP1 
## Negative:  S100A8, GZMB, NKG7, FGFBP2, CCL4, GNLY, CST7, PRF1, GZMA, S100A9 
##     LGALS2, SPON2, GZMH, CCL3, S100A12, RBP7, CTSW, MS4A6A, CCL5, CLIC3 
##     FOLR3, XCL2, GSTP1, TYROBP, TTC38, IGFBP7, ASGR1, CD160, HOPX, LYZ
\end{verbatim}

\begin{Shaded}
\begin{Highlighting}[]
\FunctionTok{ElbowPlot}\NormalTok{(pbmc)}
\end{Highlighting}
\end{Shaded}

\pandocbounded{\includegraphics[keepaspectratio]{pbmc_analysis_files/pbmc_analysis_files/figure-latex/unnamed-chunk-6-1.pdf}}

\begin{center}\rule{0.5\linewidth}{0.5pt}\end{center}

\subsection{6. Clustering}\label{clustering}

\begin{Shaded}
\begin{Highlighting}[]
\NormalTok{pbmc }\OtherTok{\textless{}{-}} \FunctionTok{FindNeighbors}\NormalTok{(pbmc, }\AttributeTok{dims =} \DecValTok{1}\SpecialCharTok{:}\DecValTok{8}\NormalTok{)}
\end{Highlighting}
\end{Shaded}

\begin{verbatim}
## Computing nearest neighbor graph
\end{verbatim}

\begin{verbatim}
## Computing SNN
\end{verbatim}

\begin{Shaded}
\begin{Highlighting}[]
\NormalTok{pbmc }\OtherTok{\textless{}{-}} \FunctionTok{FindClusters}\NormalTok{(pbmc, }\AttributeTok{resolution =} \FloatTok{0.5}\NormalTok{)}
\end{Highlighting}
\end{Shaded}

\begin{verbatim}
## Modularity Optimizer version 1.3.0 by Ludo Waltman and Nees Jan van Eck
## 
## Number of nodes: 2638
## Number of edges: 91917
## 
## Running Louvain algorithm...
## Maximum modularity in 10 random starts: 0.8743
## Number of communities: 9
## Elapsed time: 0 seconds
\end{verbatim}

\begin{center}\rule{0.5\linewidth}{0.5pt}\end{center}

\subsection{7. UMAP Visualization}\label{umap-visualization}

\begin{Shaded}
\begin{Highlighting}[]
\NormalTok{pbmc }\OtherTok{\textless{}{-}} \FunctionTok{RunUMAP}\NormalTok{(pbmc, }\AttributeTok{dims =} \DecValTok{1}\SpecialCharTok{:}\DecValTok{8}\NormalTok{)}
\end{Highlighting}
\end{Shaded}

\begin{verbatim}
## Warning: The default method for RunUMAP has changed from calling Python UMAP via reticulate to the R-native UWOT using the cosine metric
## To use Python UMAP via reticulate, set umap.method to 'umap-learn' and metric to 'correlation'
## This message will be shown once per session
\end{verbatim}

\begin{verbatim}
## 17:13:01 UMAP embedding parameters a = 0.9922 b = 1.112
\end{verbatim}

\begin{verbatim}
## 17:13:01 Read 2638 rows and found 8 numeric columns
\end{verbatim}

\begin{verbatim}
## 17:13:01 Using Annoy for neighbor search, n_neighbors = 30
\end{verbatim}

\begin{verbatim}
## 17:13:01 Building Annoy index with metric = cosine, n_trees = 50
\end{verbatim}

\begin{verbatim}
## 0%   10   20   30   40   50   60   70   80   90   100%
\end{verbatim}

\begin{verbatim}
## [----|----|----|----|----|----|----|----|----|----|
\end{verbatim}

\begin{verbatim}
## **************************************************|
## 17:13:02 Writing NN index file to temp file /var/folders/1v/6vsfq_rj6pj4q1dvdljkhtdr0000gn/T//Rtmp4u46JM/file970351a5d855
## 17:13:02 Searching Annoy index using 1 thread, search_k = 3000
## 17:13:03 Annoy recall = 100%
## 17:13:03 Commencing smooth kNN distance calibration using 1 thread with target n_neighbors = 30
## 17:13:03 Initializing from normalized Laplacian + noise (using RSpectra)
## 17:13:04 Commencing optimization for 500 epochs, with 104232 positive edges
## 17:13:04 Using rng type: pcg
## 17:13:07 Optimization finished
\end{verbatim}

\begin{Shaded}
\begin{Highlighting}[]
\FunctionTok{DimPlot}\NormalTok{(pbmc, }\AttributeTok{reduction =} \StringTok{"umap"}\NormalTok{, }\AttributeTok{label =} \ConstantTok{TRUE}\NormalTok{)}
\end{Highlighting}
\end{Shaded}

\pandocbounded{\includegraphics[keepaspectratio]{pbmc_analysis_files/pbmc_analysis_files/figure-latex/unnamed-chunk-8-1.pdf}}

\begin{center}\rule{0.5\linewidth}{0.5pt}\end{center}

\subsection{8. Marker Gene
Identification}\label{marker-gene-identification}

\begin{Shaded}
\begin{Highlighting}[]
\NormalTok{pbmc.markers }\OtherTok{\textless{}{-}} \FunctionTok{FindAllMarkers}\NormalTok{(pbmc, }\AttributeTok{only.pos =} \ConstantTok{TRUE}\NormalTok{,}
                               \AttributeTok{min.pct =} \FloatTok{0.3}\NormalTok{,}
                               \AttributeTok{logfc.threshold =} \FloatTok{0.3}\NormalTok{)}
\end{Highlighting}
\end{Shaded}

\begin{verbatim}
## Calculating cluster 0
\end{verbatim}

\begin{verbatim}
## For a (much!) faster implementation of the Wilcoxon Rank Sum Test,
## (default method for FindMarkers) please install the presto package
## --------------------------------------------
## install.packages('devtools')
## devtools::install_github('immunogenomics/presto')
## --------------------------------------------
## After installation of presto, Seurat will automatically use the more 
## efficient implementation (no further action necessary).
## This message will be shown once per session
\end{verbatim}

\begin{verbatim}
## Calculating cluster 1
\end{verbatim}

\begin{verbatim}
## Calculating cluster 2
\end{verbatim}

\begin{verbatim}
## Calculating cluster 3
\end{verbatim}

\begin{verbatim}
## Calculating cluster 4
\end{verbatim}

\begin{verbatim}
## Calculating cluster 5
\end{verbatim}

\begin{verbatim}
## Calculating cluster 6
\end{verbatim}

\begin{verbatim}
## Calculating cluster 7
\end{verbatim}

\begin{verbatim}
## Calculating cluster 8
\end{verbatim}

\begin{Shaded}
\begin{Highlighting}[]
\FunctionTok{head}\NormalTok{(pbmc.markers)}
\end{Highlighting}
\end{Shaded}

\begin{verbatim}
##               p_val avg_log2FC pct.1 pct.2     p_val_adj cluster  gene
## RPS12 4.995816e-167  0.7782779 1.000 0.991 6.851262e-163       0 RPS12
## RPS6  2.450492e-154  0.6962518 1.000 0.995 3.360605e-150       0  RPS6
## RPS27 2.076456e-151  0.7396658 0.999 0.991 2.847652e-147       0 RPS27
## RPS25 2.801125e-148  0.8157974 0.999 0.973 3.841463e-144       0 RPS25
## RPS14 7.001491e-145  0.6539509 1.000 0.994 9.601845e-141       0 RPS14
## RPL32 3.487487e-143  0.6193642 0.999 0.995 4.782740e-139       0 RPL32
\end{verbatim}

\subsubsection{Top markers per cluster}\label{top-markers-per-cluster}

\begin{Shaded}
\begin{Highlighting}[]
\NormalTok{top10 }\OtherTok{\textless{}{-}}\NormalTok{ pbmc.markers }\SpecialCharTok{\%\textgreater{}\%}
  \FunctionTok{group\_by}\NormalTok{(cluster) }\SpecialCharTok{\%\textgreater{}\%}
  \FunctionTok{top\_n}\NormalTok{(}\AttributeTok{n =} \DecValTok{10}\NormalTok{, }\AttributeTok{wt =}\NormalTok{ avg\_log2FC)}

\NormalTok{top10}
\end{Highlighting}
\end{Shaded}

\begin{verbatim}
## # A tibble: 90 x 7
## # Groups:   cluster [9]
##        p_val avg_log2FC pct.1 pct.2 p_val_adj cluster gene     
##        <dbl>      <dbl> <dbl> <dbl>     <dbl> <fct>   <chr>    
##  1 6.26e-131       1.26 0.906 0.572 8.59e-127 0       LDHB     
##  2 3.31e-102       1.20 0.858 0.37  4.54e- 98 0       CD3D     
##  3 4.25e- 79       2.35 0.409 0.102 5.83e- 75 0       CCR7     
##  4 8.22e- 64       1.37 0.638 0.336 1.13e- 59 0       NOSIP    
##  5 1.24e- 52       2.16 0.33  0.097 1.70e- 48 0       PRKCQ-AS1
##  6 3.59e- 47       2.03 0.317 0.098 4.93e- 43 0       LEF1     
##  7 2.08e- 45       1.54 0.425 0.175 2.85e- 41 0       PIK3IP1  
##  8 5.93e- 34       1.34 0.373 0.171 8.13e- 30 0       TCF7     
##  9 4.39e- 28       1.21 0.362 0.172 6.03e- 24 0       RGCC     
## 10 7.78e- 25       1.20 0.429 0.245 1.07e- 20 0       C12orf57 
## # i 80 more rows
\end{verbatim}

\begin{center}\rule{0.5\linewidth}{0.5pt}\end{center}

\subsection{9. Heatmap of Marker Genes}\label{heatmap-of-marker-genes}

\begin{Shaded}
\begin{Highlighting}[]
\FunctionTok{DoHeatmap}\NormalTok{(pbmc, }\AttributeTok{features =}\NormalTok{ top10}\SpecialCharTok{$}\NormalTok{gene) }\SpecialCharTok{+} \FunctionTok{NoLegend}\NormalTok{()}
\end{Highlighting}
\end{Shaded}

\begin{verbatim}
## Warning in DoHeatmap(pbmc, features = top10$gene): The following features were
## omitted as they were not found in the scale.data layer for the RNA assay:
## NDRG2, PLD4, BASP1, PRSS23, AKR1C3, TPPP3, MS4A4A, HMOX1, CTSL, LILRA3, HES4,
## CDKN1C, RP11-290F20.3, CD8A, TSPAN13, BANK1, FCER2, VPREB3, SUSD3, TTC39C,
## TRADD, SPOCK2, TRAT1, CD40LG, BST1, CD14, C12orf57, RGCC, TCF7, PIK3IP1, LEF1,
## PRKCQ-AS1, NOSIP, CCR7, CD3D, LDHB
\end{verbatim}

\pandocbounded{\includegraphics[keepaspectratio]{pbmc_analysis_files/pbmc_analysis_files/figure-latex/unnamed-chunk-11-1.pdf}}

\begin{center}\rule{0.5\linewidth}{0.5pt}\end{center}

\subsection{10. Cell Type Annotation}\label{cell-type-annotation}

\begin{Shaded}
\begin{Highlighting}[]
\NormalTok{new.cluster.ids }\OtherTok{\textless{}{-}} \FunctionTok{c}\NormalTok{(}\StringTok{"Naive CD4 T"}\NormalTok{, }\StringTok{"CD14+ Mono"}\NormalTok{, }\StringTok{"Memory CD4 T"}\NormalTok{,}
                     \StringTok{"B"}\NormalTok{, }\StringTok{"CD8 T"}\NormalTok{, }\StringTok{"FCGR3A+ Mono"}\NormalTok{, }\StringTok{"NK"}\NormalTok{, }\StringTok{"DC"}\NormalTok{, }\StringTok{"Platelet"}\NormalTok{)}

\FunctionTok{names}\NormalTok{(new.cluster.ids) }\OtherTok{\textless{}{-}} \FunctionTok{levels}\NormalTok{(pbmc)}

\NormalTok{pbmc }\OtherTok{\textless{}{-}} \FunctionTok{RenameIdents}\NormalTok{(pbmc, new.cluster.ids)}
\end{Highlighting}
\end{Shaded}

\subsubsection{UMAP with cell type
labels}\label{umap-with-cell-type-labels}

\begin{Shaded}
\begin{Highlighting}[]
\FunctionTok{DimPlot}\NormalTok{(pbmc, }\AttributeTok{reduction =} \StringTok{"umap"}\NormalTok{, }\AttributeTok{label =} \ConstantTok{TRUE}\NormalTok{, }\AttributeTok{pt.size =} \FloatTok{0.5}\NormalTok{)}
\end{Highlighting}
\end{Shaded}

\pandocbounded{\includegraphics[keepaspectratio]{pbmc_analysis_files/pbmc_analysis_files/figure-latex/unnamed-chunk-13-1.pdf}}

\begin{center}\rule{0.5\linewidth}{0.5pt}\end{center}

\subsection{11. Differential Expression
Analysis}\label{differential-expression-analysis}

\begin{Shaded}
\begin{Highlighting}[]
\NormalTok{mono.de }\OtherTok{\textless{}{-}} \FunctionTok{FindMarkers}\NormalTok{(pbmc, }\AttributeTok{ident.1 =} \StringTok{"CD14+ Mono"}\NormalTok{, }\AttributeTok{min.pct =} \FloatTok{0.25}\NormalTok{)}

\FunctionTok{head}\NormalTok{(mono.de, }\DecValTok{10}\NormalTok{)}
\end{Highlighting}
\end{Shaded}

\begin{verbatim}
##                p_val avg_log2FC pct.1 pct.2     p_val_adj
## S100A8  0.000000e+00   6.338298 0.984 0.137  0.000000e+00
## LGALS2  0.000000e+00   4.999301 0.899 0.079  0.000000e+00
## S100A9 1.196360e-312   5.835831 0.995 0.232 1.640688e-308
## FCN1   9.511239e-284   3.662344 0.952 0.167 1.304371e-279
## CD14   1.204764e-268   5.173508 0.671 0.040 1.652213e-264
## TYROBP 1.079389e-254   3.092964 0.993 0.280 1.480274e-250
## MS4A6A 3.374271e-251   4.895105 0.687 0.054 4.627476e-247
## LYZ    2.299040e-241   4.292017 1.000 0.527 3.152904e-237
## CST3   5.613179e-234   2.898334 0.991 0.281 7.697914e-230
## TYMP   1.033394e-206   2.865588 0.910 0.238 1.417197e-202
\end{verbatim}

\begin{center}\rule{0.5\linewidth}{0.5pt}\end{center}

\subsection{12. Key Findings}\label{key-findings}

\begin{itemize}
\tightlist
\item
  \textasciitilde2600 cells retained after quality control filtering
\item
  9 distinct immune cell populations identified
\item
  Clear separation of clusters in UMAP space
\item
  Marker genes align with known PBMC biology
\end{itemize}

\begin{center}\rule{0.5\linewidth}{0.5pt}\end{center}

\subsection{13. Session Info}\label{session-info}

\begin{Shaded}
\begin{Highlighting}[]
\FunctionTok{sessionInfo}\NormalTok{()}
\end{Highlighting}
\end{Shaded}

\begin{verbatim}
## R version 4.5.3 (2026-03-11)
## Platform: x86_64-apple-darwin20
## Running under: macOS Sequoia 15.7.4
## 
## Matrix products: default
## BLAS:   /Library/Frameworks/R.framework/Versions/4.5-x86_64/Resources/lib/libRblas.0.dylib 
## LAPACK: /Library/Frameworks/R.framework/Versions/4.5-x86_64/Resources/lib/libRlapack.dylib;  LAPACK version 3.12.1
## 
## locale:
## [1] en_US.UTF-8/en_US.UTF-8/en_US.UTF-8/C/en_US.UTF-8/en_US.UTF-8
## 
## time zone: America/Toronto
## tzcode source: internal
## 
## attached base packages:
## [1] stats     graphics  grDevices utils     datasets  methods   base     
## 
## other attached packages:
##  [1] future_1.70.0      cowplot_1.2.0      lubridate_1.9.4    forcats_1.0.0     
##  [5] stringr_1.5.1      dplyr_1.1.4        purrr_1.0.4        readr_2.1.5       
##  [9] tidyr_1.3.1        tibble_3.2.1       ggplot2_4.0.2      tidyverse_2.0.0   
## [13] Seurat_5.4.0       SeuratObject_5.4.0 sp_2.2-1          
## 
## loaded via a namespace (and not attached):
##   [1] deldir_2.0-4           pbapply_1.7-4          gridExtra_2.3         
##   [4] rlang_1.1.6            magrittr_2.0.3         RcppAnnoy_0.0.23      
##   [7] spatstat.geom_3.7-3    matrixStats_1.5.0      ggridges_0.5.7        
##  [10] compiler_4.5.3         png_0.1-9              vctrs_0.6.5           
##  [13] reshape2_1.4.5         pkgconfig_2.0.3        fastmap_1.2.0         
##  [16] labeling_0.4.3         utf8_1.2.4             promises_1.3.3        
##  [19] rmarkdown_2.31         tzdb_0.5.0             xfun_0.52             
##  [22] jsonlite_2.0.0         goftest_1.2-3          later_1.4.2           
##  [25] spatstat.utils_3.2-2   irlba_2.3.7            parallel_4.5.3        
##  [28] cluster_2.1.8.2        R6_2.6.1               ica_1.0-3             
##  [31] stringi_1.8.7          RColorBrewer_1.1-3     spatstat.data_3.1-9   
##  [34] reticulate_1.46.0      parallelly_1.46.1      spatstat.univar_3.1-7 
##  [37] lmtest_0.9-40          scattermore_1.2        Rcpp_1.1.0            
##  [40] knitr_1.50             tensor_1.5.1           future.apply_1.20.2   
##  [43] zoo_1.8-15             R.utils_2.13.0         sctransform_0.4.3     
##  [46] timechange_0.3.0       httpuv_1.6.16          Matrix_1.7-4          
##  [49] splines_4.5.3          igraph_2.2.3           tidyselect_1.2.1      
##  [52] rstudioapi_0.17.1      abind_1.4-8            yaml_2.3.10           
##  [55] spatstat.random_3.4-5  codetools_0.2-20       miniUI_0.1.2          
##  [58] spatstat.explore_3.8-0 listenv_0.10.1         lattice_0.22-9        
##  [61] plyr_1.8.9             withr_3.0.2            shiny_1.11.1          
##  [64] S7_0.2.1               ROCR_1.0-12            evaluate_1.0.3        
##  [67] Rtsne_0.17             fastDummies_1.7.5      survival_3.8-6        
##  [70] polyclip_1.10-7        fitdistrplus_1.2-6     pillar_1.10.2         
##  [73] KernSmooth_2.23-26     plotly_4.12.0          generics_0.1.3        
##  [76] RcppHNSW_0.6.0         hms_1.1.3              scales_1.4.0          
##  [79] globals_0.19.1         xtable_1.8-4           glue_1.8.0            
##  [82] lazyeval_0.2.3         tools_4.5.3            data.table_1.17.0     
##  [85] RSpectra_0.16-2        RANN_2.6.2             dotCall64_1.2         
##  [88] grid_4.5.3             nlme_3.1-168           patchwork_1.3.2       
##  [91] cli_3.6.5              spatstat.sparse_3.1-0  spam_2.11-3           
##  [94] viridisLite_0.4.2      uwot_0.2.4             gtable_0.3.6          
##  [97] R.methodsS3_1.8.2      digest_0.6.37          progressr_0.19.0      
## [100] ggrepel_0.9.8          htmlwidgets_1.6.4      farver_2.1.2          
## [103] R.oo_1.27.1            htmltools_0.5.8.1      lifecycle_1.0.4       
## [106] httr_1.4.7             mime_0.13              MASS_7.3-65
\end{verbatim}

\end{document}
